%% =============================================================================
%% ANEXO: MANUAL DE DESPLIEGUE
%% =============================================================================
\chapter{Manual de despliegue}
\label{apx:despliegue}

Este anexo describe el procedimiento para desplegar el sistema en un entorno de
desarrollo. El repositorio se organiza como un monorepo cuya estructura de alto
nivel se muestra a continuación.

\begin{dirtreebox}[Estructura del repositorio]
  \dirtreeitem[0]{tfg-robot-patrulla/}
  \dirtreeitem[1]{backend/}
  \dirtreeitem[1]{flutter\_frontend/}
  \dirtreeitem[1]{robot\_ws/}
  \dirtreeitem[1]{scripts/}
  \dirtreeitem[1]{docs/}
\end{dirtreebox}

Cada directorio agrupa, respectivamente, el backend y sus consumidores de
eventos (\texttt{backend}), la aplicación de usuario (\texttt{flutter\_frontend}),
el \textit{workspace} de \gls{ros} con los nodos del robot (\texttt{robot\_ws}),
los \textit{scripts} de automatización del despliegue (\texttt{scripts}) y la
documentación, incluida esta memoria (\texttt{docs}).

%% =============================================================================
\section{Requisitos previos}
\label{sec:requisitos-despliegue}

El despliegue requiere un sistema con Docker y Docker Compose para la
infraestructura de datos, Python para el backend, el SDK de Flutter para la
aplicación de usuario y una instalación de \gls{ros} Humble para la capa
robótica. El conjunto de \textit{scripts} incluido en el directorio
\texttt{scripts/} automatiza la comprobación de prerrequisitos y el arranque de
cada componente.

%% =============================================================================
\section{Procedimiento de despliegue}
\label{sec:procedimiento-despliegue}

El sistema se levanta por capas, respetando el orden de dependencias. En primer
lugar se comprueban los prerrequisitos y se arranca la infraestructura de datos
(bus de mensajería, bases de datos y almacenamiento de objetos); a continuación
se preparan y ejecutan el backend y la aplicación de usuario; por último, se
construye y se lanza la capa robótica.

\begin{bashcode}[title={Despliegue de la infraestructura y el backend}]
./scripts/00_prereqs.sh        # Comprobación de prerrequisitos
./scripts/10_infra_up.sh       # Infraestructura de datos (Docker Compose)
./scripts/01_backend_setup.sh  # Preparación del entorno del backend
./scripts/11_seed.sh           # Carga de datos iniciales
./scripts/20_backend_run.sh    # Arranque del backend
\end{bashcode}

Una vez disponibles la infraestructura y el backend, se arranca la aplicación de
usuario y, finalmente, la capa robótica sobre el robot.

\begin{bashcode}[title={Despliegue del frontend y de la capa robótica}]
./scripts/30_frontend_run.sh   # Aplicación de usuario (Flutter)
./scripts/41_robot_build.sh    # Construcción del workspace ROS 2
./scripts/42_robot_run.sh      # Arranque de los nodos del robot
\end{bashcode}

%% =============================================================================
\section{Configuración mediante variables de entorno}
\label{sec:configuracion-despliegue}

Los parámetros sensibles y dependientes del entorno (credenciales de las bases de
datos, direcciones del bus de mensajería, identificador del robot y claves de
acceso al almacenamiento de objetos) se gestionan mediante variables de entorno,
a partir de las plantillas \texttt{.env.example} incluidas en el repositorio.
Antes del primer arranque debe crearse el fichero \texttt{.env} correspondiente y
completarse con los valores adecuados.
